
\chapter{Introduzione}

Da un punto di vista ingegneristico, l'esplorazione subacquea ha storicamente rivestito un ruolo di primaria importanza, in virtù delle molteplici possibili applicazioni in ambito commerciale, scientifico e militare. Tradizionalmente, però, la necessità di garantire la presenza di personale umano a bordo ha sempre rappresentato un fattore fortemente limitante, ostacolando il pieno sfruttamento delle potenzialità offerte da questo settore. Nel corso degli ultimi decenni, tuttavia, i notevoli avanzamenti tecnologici nel campo della robotica e della sensoristica hanno portato ad un cambio di paradigma. Si sono, infatti, diffusi i cosiddetti \emph{Unmanned Underwater Vehicles (UUV)}, veicoli subacquei in grado di opeare in assenza di equipaggio, seppur con vari livelli di autonomia. All'interno della definizione di UUV rientrano veicoli anche molto differenti tra loro, sia per caratteristiche che per area di impiego. In particolare, questi possono essere racchiusi in due macrocategorie: gli \emph{Autonomous Underwater Vehicles (AUV)} e i \emph{Remotely Operated Underwater Vehicles}. In questa trattazione viene considerata solo la prima classe di droni, essendo più attinente agli scopi del progetto. Come facilmente intuibile, essi sono caratterizzati da una completa autonomia durante ogni fase della navigazione. Questa caratteristica li rende la scelta più indicata per missioni di esplorazione, rilevamento e raccolta dati su lunghe distanze o in luoghi difficilmente raggiungibili dall'uomo. Alcuni utilizzi tipici sono rappresentati da missioni di mappatura del fondale, di studio della fauna marittima, di localizzazione e ispezione di gasdotti e di analisi delle principali proprietà dell'acqua marina (come temperatura e salinità). 

Lo sviluppo di veicoli di questo tipo pone importanti sfide ingegneristiche, legate non solo alla loro progettazione e realizzazione ma anche alle condizioni particolarmente ostili dell'ambiente subacqueo. A causa dell'elevata densità dell'acqua e delle sue caratteristiche, infatti, i principali sistemi di posizionamento esterni, basati sul segnale satellitare, non possono essere adottati se non nelle fasi di navigazione superficiale. Sebbene l'impiego di altre tecniche sia possibile, esse risultano non sempre concretamente utilizzabili e spesso economicamente non vantaggiose. Una delle principali soluzioni consiste nell'uso del \emph{Long Baseline System (LBL)}, capace di determinare la posizione del drone relativamente ad alcuni transponder posti sul fondale attraverso l'invio e la ricezione di un segnale acustico. Una possibile variante di tale sistema è lo \emph{Short Baseline System (SBL)}, basato su array di transpoder ridotti e, pertanto, più indicato per missioni in acque poco profonde o in cui è richiesta un'accuratezza inferiore. Un'alternativa simile, denominata \emph{Beacon-based Localization}, prevede, invece, l'utilizzo di fari acustici o elettromagnetici posti in precise località. I metodi elencati, tuttavia, dipendono dalla presenza di oggetti esterni, il cui posizionamento può essere complesso se non addirittura impossibile. 

Solitamente, dunque, è necessario fare affidamento sui soli sensori interni al drone, che possono però fornire una stima della posizione unicamente mediante un processo di regressione. Tra le varie possibilità, il sensore in assoluto più diffuso risulta essere l'\emph{Inertial Measurement Unit (IMU)}, facente parte dell'\emph{Inertial Navigation System (INS)} e capace di determinare l'assetto, le componenti di velocità angolare e le accelerazioni caratteristiche del drone grazie all'utilizzo combinato di un giroscopio e di un accelerometro. Un simile sensore, sebbene sia molto versatile e consenta di valutare l'evoluzione della posizione nel tempo tramite una semplice integrazione, tende ad essere molto sensibile. Il processo costruttivo del sensore stesso e il montaggio sull'AUV, infatti, introducono inevitabilmente degli errori di misurazione. Nonostante una ricalibrazione periodica possa incrementare l'accuratezza della stima, la deriva rispetto alla posizione reale, specialmente per grandi distanze, rimane significativa. Solitamente, dunque, si sceglie di affiancarlo ad un \emph{Doppler Velocity Log (DVL)}, in grado di sfruttare proprio l'effetto Doppler per fornire una stima del vettore velocità dell'AUV. Nel dettaglio, quattro raggi acustici, la cui inclinazione rispetto agli assi body del drone è nota, vengono inviati in direzione del fondale. In base al tempo trascorso prima della ricezione è possibile stimare, quindi, la velocità di ognuno di essi rispetto al veicolo e, mediante l'utilizzo di un'apposita matrice di trasformazione, di giungere alle componenti di velocità del drone stesso. Data l'elevata accuratezza che caratterizza tale sensore, le sue misurazioni sono usate per correggere periodicamente le stime prodotte dal solo INS; tuttavia, non è esente da errori di misura e presenta una bassa frequenza di aggiornamento. Inoltre, specialmente in presenza di un fondale poco regolare, uno o più raggi possono essere deflessi in direzioni incompatibili con la ricezione. In questo caso, la regressione del vettore velocità perde accuratezza e può persino risultare impossibile. Infine, per quanto concerne la stima della profondità del drone si tende ad utilizzare degli \emph{Hydrostatic Pressure Sensors (HPS)}, solitamente molto affidabili e precisi. In particolare, essi stimano la quota dell'AUV sulla base della pressione dell'acqua circostante.

Come accennato, il posizionamento del drone avviene mediante un processo di \emph{data fusion}. Tradizionalmente ciò richiede l'utilizzo di un algoritmo model-based, che per definizione si basa sulla creazione di un modello del drone nel suo complesso e dell'ambiente circostante. La tecnica più diffusa in questo ambito risulta senza dubbio essere il \emph{filtro di Kalman}. Nonostante esso possa filtrare le informazioni provenienti dai sensori per ridurre l'impatto degli errori di misura costanti sulla regressione della posizione, difficilmente un modello risulta essere privo di approssimazioni. Queste, naturalmente, portano a discrepanze rispetto al sistema reale, introducendo una nuova fonte di errore. 

Questa tesi nasce, nell'ambito di una collaborazione con l'azienda Modelway, proprio per fronteggiare queste limitazioni. Negli ultimi anni, infatti, il rapido incremento della potenza computazionale disponibile ha favorito la ricerca su approcci alternativi, in grado di sfruttare i vantaggi derivanti da una fusione delle misurazioni provenienti dai differenti sensori senza necessariamente aver bisogno di un modello del sistema. Simili approcci, definiti data-driven poiché necessitano unicamente della conoscenza dei dati, sono ascrivibili all'ambito dell'intelligenza artificiale e, in particolare, del deep learning. 

Nel dettaglio, la tesi risulta complessivamente suddivisa in 7 capitoli. Nel Capitolo 2 viene presentata una dettagliata panoramica sull'intelligenza artificiale, esaminando in particolare il meccanismo di funzionamento e i principali elementi di una rete neurale. Nel Capitolo 3 sono analizzate le principali architetture neurali proposte in letteratura per l'ottimizzazione della regressione della posizione di un AUV. Nel Capitolo 4 è contenuta una descrizione del procedimento seguito per generare un dataset rappresentativo delle tipiche missioni che un drone commerciale autonomo potrebbe dover svolgere nella sua vita operativa. Nel Capitolo 5 si illustrano, invece, i risultati dell'analisi statistica di correlazione svolta al fine di ridurre la dimensionalità del dataset, individuando ed eliminando le misurazioni non significative per la rete neurale. Nel Capitolo 6 sono discussi i risultati dell'addestramento su alcune varianti di NavNet, l'architettura neurale selezionata per questa tesi. Infine, il Capitolo 7 riporta le conclusioni a cui si è giunti, con un'indicazione su possibili sviluppi futuri.  